\section{Outward propagation and the excess-charge bound}
\label{m:sec:outward}

The inverse comparisons now allow us to replace a reaction term by the
conditional electron density on one more layer of local scales.  We perform
this replacement before forgetting the observations at the current scale.
The failed comparisons contribute a nonnegative error density.  Its expected
integral, rather than the probability of a simultaneous comparison throughout
space, is the quantity that we shall sum.

We use the scales $r_j=2^j r_0$, $r_n=r_*$ and the retained data
$\mathcal F_j$ of Section~\ref{m:sec:conditioning}.
The data $\mathcal F_n$ are trivial.  Write $\mu_j$ for their conditional
smoothed density. At each step we take a local maximum of the previous
comparison field and the observed field $H_j$, after lowering them by
different constants. The previous field supplies the nonlinear reaction;
$H_j$ supplies the actual density through its Poisson equation. The unequal
shifts let the inverse comparisons prove the missing source inequality
for each branch that can attain the maximum.

After the last observations are forgotten, the resulting positive exterior
field will bound the smoothed electron mass in $\{a<r_*\}$ by $Z$ plus
the accumulated error. The exterior-count estimate of
Lemma~\ref{m:lem:void-count} controls the remaining smoothed mass, while
the error has integral $o(M)$. These are the two contributions to the
final excess-charge estimate.

The invariant at scale $r=r_j$ is a continuous locally
$H^1$ offset $u$ and a nonnegative bounded compactly supported error density
$d_j$ such that
\begin{equation}
\label{m:eq:outward-invariant}
 \Delta u\ge -4\pi\nu_s
       +4\pi\1_{\{a<r\}}\mu_j-4\pi d_j
       +\1_{\{a\ge r\}}f_y(u)
 \quad\text{in distributions},
\end{equation}
and
\begin{equation}
\label{m:eq:outward-bounds}
 b_r\le u\le C_{\mathrm{inv}}a^{-4}\quad\text{on }\{a\ge r\},
 \qquad u=b_r\quad\text{on }\{a>L_1r\}.
\end{equation}
Here and below an inequality between spatial functions is asserted for every
positive-density retained path, using the versions of
Lemma~\ref{lem:obs-data}.  The functions $u,d_j$ are jointly measurable
in the retained data and space.  We also maintain deterministic bounds, for the
fixed physical system and step, on the modulus of continuity and the $H^1$
norm of $u$ on every compact set, and on the size and support of $d_j$.
These additional bounds justify conditional integration.  The integral
estimate for $d_j$ will instead have universal constants.

\subsection{Replacing one layer of reaction by density}

For this subsection put $r=r_j$, $R_+=2r$, and $H=H_j$.  All parameters are fixed in
the order specified in Section~\ref{m:sec:conditioning}.  We record the
inequalities that enter the argument:
\begin{equation}
\label{m:eq:outward-shifts}
 \lambda_1<\frac\gamma4,\qquad
 16\delta_0+2e_0<\lambda_1-\lambda_2,\qquad
 \delta_0<\lambda_2,\qquad
 2e_0<c'(2L_1)^{-4}.
\end{equation}
The constant $c'>0$ is the margin in $b_r\ge h_b+c'a^{-4}$.
Set $\eta=e_0R_+^{-4}$ and $T_y=f_y/(4\pi)$.

For the current data, denote the high and low failure sets of
Lemma~\ref{m:lem:cond-inverse} by
\begin{align*}
 \mathcal B^{\mathrm h}
  &=\{y:\mu_j(y)\ge T_y(h_b(y)),\quad
        H(y)<I_y(\mu_j(y))-\delta_0a_y^{-4}\},\\
 \mathcal B^{\mathrm l}
  &=\{y:\mu_j(y)\le T_y(h_{\mathrm{top}}(y)),\quad
        H(y)>I_y(\mu_j(y))+\delta_0a_y^{-4}\}.
\end{align*}
These sets are used only within the participating band
$r\le a\le2L_1R_+$.  Let $\Pi$, $B_Q$, $Q'$ and $\mathcal E$ be the penalty,
its cube coefficients, the enlarged cubes and their bad union from
Lemma~\ref{m:lem:cond-cap}.  In particular $\Pi=0$ off $\mathcal E$ and
$H-\Pi\le C_{\mathrm{cap}}a^{-4}$ throughout this band.

\begin{proposition}[One outward step]
\label{m:prop:outward-step}
Suppose that $u,d_j$ satisfy the inductive conditions above, including
\eqref{m:eq:outward-invariant}--\eqref{m:eq:outward-bounds}, at scale $r=r_j$.
There are an offset $u'$ and a nonnegative error density $\widetilde d$,
measurable with respect to the same data $\mathcal F_j$, that satisfy these
inequalities at scale $R_+$ with $\mu_j$ still in the core term.  One may take
\begin{align}
\label{m:eq:outward-error}
 \widetilde d={}&d_j+C r^{-2}\sum_Q B_Q\1_{Q'}\notag\\
 &+\1_{\{r\le a<R_+\}}\mu_j
                       \1_{\mathcal E\cup\mathcal B^{\mathrm h}}\notag\\
 &+\1_{\{R_+\le a\le2L_1R_+\}}T_y(H)
                       \1_{\mathcal E\cup\mathcal B^{\mathrm l}},
\end{align}
where $C$ is a universal constant large enough to dominate the penalty
Laplacian.
\end{proposition}

\begin{proof}
Define two shifted offsets
\[
 u_1=u-\lambda_1R_+^{-4},\qquad
 u_2=H-\lambda_2R_+^{-4}-\Pi.
\]
Use the following open cover to define $u'$:
\begin{equation}
\label{m:eq:outward-branches}
\begin{array}{c|c}
 \text{open region}&u'\\ \hline
 a<(1+\kappa)R_+&\max(u_1,u_2)\\
 (1+\kappa/2)R_+<a<2L_1R_+&\max(u_1,u_2,b_{R_+})\\
 a>L_1R_+&b_{R_+}.
\end{array}
\end{equation}
The definitions agree on overlaps.  On the lower overlap,
Lemma~\ref{m:lem:barrier-gap} and \eqref{m:eq:outward-shifts} give
\[
 u_1\ge b_r-\lambda_1R_+^{-4}>b_{R_+}.
\]
On the upper overlap the shifted cap inequalities, ensured by the choice of
$L_1$, put both $u_1$ and $u_2$ below $b_{R_+}$.  Thus the local maxima glue
to a continuous locally $H^1$ function.  No assertion about the measure or
regularity of a level set of $a$ is required.  The new exterior bounds follow
from the same comparisons: for $R_+\le a<(1+\kappa)R_+$ the lower bound is
supplied by $u_1$, and elsewhere the barrier is an included branch.

We verify the differential inequality locally.  At a fixed point in one of
the open regions, discard branches more than $\eta$ below the maximum.
There is a neighborhood on which these discarded branches cannot attain
the maximum, while every retained branch is within $2\eta$ of it.  We show
that each retained branch satisfies the common inequality
\begin{equation}
\label{m:eq:outward-common}
 \Delta v\ge -4\pi\nu_s+4\pi\1_{\{a<R_+\}}\mu_j
             -4\pi\widetilde d+\1_{\{a\ge R_+\}}f_y(v)
\end{equation}
on that neighborhood.

For $u_1$, the old core inequality suffices on $a<r$, and monotonicity of
$f_y$ gives the required reaction on $a\ge R_+$.  On the gain region
$r\le a<R_+$, the error in \eqref{m:eq:outward-error} pays $\mu_j$ at points
of $\mathcal E\cup\mathcal B^{\mathrm h}$.  At every other gain point,
$\Pi=0$ and near activity yields
\[
 H\le u-(\lambda_1-\lambda_2)R_+^{-4}+2\eta.
\]
Here $h_b\le u\le h_{\mathrm{top}}$.  If $\mu_j>T_y(u)$, then the high
comparison applies and $I_y(\mu_j)\ge u$, including the case where the
inverse is clipped at its upper endpoint.  It would follow that
\[
 H\ge u-\delta_0a^{-4}\ge u-16\delta_0R_+^{-4},
\]
contrary to \eqref{m:eq:outward-shifts}.  Thus $4\pi\mu_j\le f_y(u)$.
The old reaction supplies precisely the new core term on the gain region.

For $u_2$, the identity $\Delta H=-4\pi\nu_s+4\pi\mu_j$ gives the required
core term on all of $a<R_+$, while the first increment in
\eqref{m:eq:outward-error} pays the Laplacian of $\Pi$.  Consider the remaining
points where $u_2$ is retained.  The lower splice or the included barrier gives
\[
 u_2\ge b_{R_+}-2\eta
       \ge h_b+c'a^{-4}-2e_0R_+^{-4}>h_b.
\]
The last inequality uses $a\le2L_1R_+$ and
\eqref{m:eq:outward-shifts}.  On
$\mathcal E\cup\mathcal B^{\mathrm l}$ the last error term pays the
reaction because $u_2\le H$ and $f_y$ is nondecreasing.  Otherwise $\Pi=0$
and the cap gives $H<h_{\mathrm{top}}$.  If $\mu_j<T_y(u_2)$, the low
comparison applies and $I_y(\mu_j)\le u_2$, also when the inverse is clipped
at its lower endpoint.  We would obtain
\[
 H\le u_2+\delta_0a^{-4}\le u_2+\delta_0R_+^{-4},
\]
contradicting $H=u_2+\lambda_2R_+^{-4}$.  Therefore
$4\pi\mu_j\ge f_y(u_2)$, as required.

The barrier branch appears only where $a>R_+$, and
Lemma~\ref{m:lem:barrier-subsolution} gives \eqref{m:eq:outward-common} there.
The old error density and every additional error are nonnegative, so they
can only weaken the required lower bound.

These comparisons cover the exact level sets as well: $a=r$ belongs to the
gain region and uses the old reaction, whereas $a=R_+$ belongs to the new
exterior and uses its reaction.  All right-side comparisons hold almost
everywhere on the chosen neighborhood and therefore upgrade each retained
branch's distributional inequality to \eqref{m:eq:outward-common}.
The source $\nu_s$ is a smooth bounded density for the fixed system,
and all branches here are continuous and locally $H^1$. The reaction
is also bounded on bounded ranges of its argument: for a fixed system,
$t_{\min}:=\inf_z t_z>0$ and $V_c\in L_c^{5/2}$ give, for every $A>0$,
\[
 \sup_{y,\,|h|\le A} f_y(h)
 \le C t_{\min}^{-3}\int V_c^{3/2}
       +CA^{3/2}\sup_y R1(y)<\infty.
\]
Thus the first,
nonsingular part of Lemma~\ref{lem:maximum}, followed by a partition
of unity, proves the inequality for $u'$ on all of $\R^3$. No vanishing
of the reaction near the original nuclei is needed in this step.
\end{proof}

\subsection{Averaging the data and summing the error}

The next step explains why the nonlinear reaction survives the removal of
observations.  It also records the integrability needed in that operation.

\begin{lemma}[Forgetting the current observations]
\label{m:lem:outward-average}
Set
\[
 u_{j+1}=\E[u'\mid\mathcal F_{j+1}],\qquad
 d_{j+1}=\E[\widetilde d\mid\mathcal F_{j+1}].
\]
Then $u_{j+1},d_{j+1}$ satisfy
\eqref{m:eq:outward-invariant}--\eqref{m:eq:outward-bounds} at scale $r_{j+1}$,
with conditional density $\mu_{j+1}$.  They retain continuity and local
$H^1$ regularity for the offset and bounded compact support for the
nonnegative error density.
\end{lemma}

\begin{proof}
For the fixed system there are finitely many scales and cubes, and all
master smoothing widths have a positive minimum.
Lemma~\ref{lem:obs-data} therefore gives deterministic bounds for the
conditional densities, their potentials and their gradients.  The offsets
and error densities in the assertion are jointly measurable: maxima,
clipped inverses and failure flags preserve measurability, and each cube
supremum may be taken over a fixed countable dense subset.

Inductively the branches have deterministic continuity moduli and local
$H^1$ bounds on each compact set.  Their open definitions agree on overlaps;
a finite subcover of a compact set supplies a Lebesgue number and hence a
modulus for the glued function.  The gradient of a finite maximum is almost
everywhere one of the branch gradients.  Conditional averaging consequently
preserves continuity.  It preserves the local $H^1$ bound by testing weak
derivatives against compactly supported vector fields, using Fubini and
Cauchy--Schwarz.  The error terms have deterministic bounded size and support
for the fixed system and step, so the same is true after averaging.

Integrate the distributional inequality of
Proposition~\ref{m:prop:outward-step} against a fixed nonnegative compact test
and then against the explicit positive-parent conditional law of
Lemma~\ref{lem:obs-data}.  The preceding bounds justify Fubini.  The tower
property gives
$\E[\mu_j\mid\mathcal F_{j+1}]=\mu_{j+1}$, while convexity gives
\[
 \E[f_y(u')\mid\mathcal F_{j+1}]
     \ge f_y\bigl(\E[u'\mid\mathcal F_{j+1}]\bigr).
\]
This is the required lower bound for the reaction.  The exterior bounds and
the exterior equality are deterministic and pass to the average.
\end{proof}

\begin{lemma}[Integrated error]
\label{m:lem:outward-error}
Starting with $u=b_{r_0}$ and $d_0=0$, the preceding construction reaches
scale $r_n=r_*$ and produces a deterministic error density with
\begin{equation}
\label{m:eq:outward-total-error}
 \int d_n\le C M r_*^{17-3w}=o(M)
 \qquad\text{along the contradiction sequence}.
\end{equation}
The constant is independent of the number of scales and of the nuclear
configuration.
\end{lemma}

\begin{proof}
Since $r_0<Z^{-1/3}\le a$, Lemma~\ref{m:lem:barrier-subsolution} gives the
initial invariant everywhere; there is no initial core region.

At scale $r=r_j$, apply Lemma~\ref{m:lem:cond-moments} to the three
increments in \eqref{m:eq:outward-error}.  The high-failure estimate pays
$\mu_j$ on the gain region; the low-failure estimate pays
$T_y(H)=f_y(H)/(4\pi)$ on the remaining participating region.
The same Lemma pays membership in $\mathcal E$ and the penalty
Laplacian. For $\mathcal B\in\{\mathcal B^{\mathrm h},\mathcal B^{\mathrm l}\}$,
use $\1_{\mathcal E\cup\mathcal B}\le\1_{\mathcal E}+\1_{\mathcal B}$.
All these integrands are nonnegative, so restriction to
the indicated subregions only improves their bounds.  We obtain
\[
 \E\int(\widetilde d-d_j)\le C M r^{17-3w}.
\]

Conditional averaging preserves the unconditional expected integral of a
nonnegative density.  Because $\beta=17-3w>0$,
\[
 \sum_{j=0}^{n-1}r_j^\beta
   =r_*^\beta\sum_{h=1}^{n}2^{-h\beta}
   \le\frac{r_*^\beta}{2^\beta-1}.
\]
At the last step the data are trivial, so $d_n$ is deterministic.  This
proves \eqref{m:eq:outward-total-error}, uniformly in $n$.  The argument has
used the failure probability separately at each spatial point and then
integrated it; it has imposed no simultaneous good-data event in space.
\end{proof}

\subsection{The total-charge comparison}

We have reached the unconditional density $\mu_n$ and replaced the reaction
by this density throughout $\{a<r_*\}$, at total error $o(M)$.  A final
potential comparison converts this statement to a bound on the number of
electrons in that region.

\begin{proof}[Proof of Theorem~\ref{thm:main}]
Use the contradiction sequence and sector eigenstates of
Section~\ref{m:sec:conditioning}. Its terminal scale is
$r_*=(B_0M/K_{\mathrm{exc}})^{1/3}$. The universal $B_0$ will exceed
the fixed count constant below.  All deterministic parameters and tolerances
are chosen before taking the sequence limit.  The smallness requirements of
the preceding lemmas then hold for all its sufficiently late members,
uniformly over their scales.

At the final step put
\[
 U=V_c+u_n,\qquad
 S=\nu-\1_{\{a<r_*\}}\mu_n+d_n.
\]
The signed measure $S$ has compact support.  In fact $\{a<r_*\}$ lies in a
finite union of balls around the nuclei, the error density has compact
support, and the nuclei are finite.  Since
$\Delta V_c=-4\pi\nu+4\pi\nu_s$, the final invariant implies
\[
 \Delta(U-K*S)\ge\1_{\{a\ge r_*\}}f_y(u_n)\ge0.
\]
The nuclear singularities cancel in the explicit identity
\begin{equation}
\label{m:eq:outward-cancellation}
 U-K*S
   =u_n-V_s+K*(\1_{\{a<r_*\}}\mu_n)-K*d_n.
\end{equation}
The right-hand side is continuous and locally $H^1$ on all of $\R^3$:
the last two potentials have bounded compactly supported densities, and the
other terms have the stated regularity from the construction.

Far from the nuclei, \eqref{m:eq:outward-bounds} gives
$U=B_{r_*}=W_0+q_0-\Gamma_{r_*}$, which tends to zero by the deterministic
field estimates.  The compact source potential $K*S$ also tends to zero.
For every $\epsilon>0$, the function $(U-K*S-\epsilon)_+$ consequently has
compact support and is an admissible $H^1$ test in the subharmonic
inequality.  Testing gives a nonpositive squared gradient norm, so this
positive part vanishes.  Letting $\epsilon\downarrow0$ proves
\[
 U\le K*S\qquad\text{away from the nuclear points}.
\]
On every sufficiently large sphere $U=B_{r_*}\ge .9W_0>0$.
The spherical mean of the potential of a compactly supported signed measure,
on a sphere centered at the origin and enclosing its support, equals its
total mass divided by the sphere radius.  Hence
\begin{equation}
\label{m:eq:outward-mass}
 0<\int\dd S
   =Z-N+\int_{\{a\ge r_*\}}\mu_n+\int d_n.
\end{equation}
Only positivity on a far sphere is used here; no asymptotic charge formula
for $W_0$ is needed.

At the last scale the data have been completely averaged away, so
\[
 \mu_n(y)=\E_\Psi\sum_i\mathcal K_{x_i}(y).
\]
If a kernel centered at $x$ contributes to a point with $a_y\ge r_*$,
then its support and the Lipschitz scale estimate imply
\[
 a_y\le a_x+\frac{|x-y|}{L}
       \le a_x\left(1+\frac{c_1r_*^w}{L}\right).
\]
For small $r_*$ this forces $a_x\ge r_*/2$.  Each kernel has integral one,
and Lemma~\ref{m:lem:void-count} therefore gives
\[
 \int_{\{a\ge r_*\}}\mu_n
   \le\E_\Psi\#\{i:a(x_i)\ge r_*/2\}
   \le C M r_*^{-3}.
\]
Combining this with \eqref{m:eq:outward-total-error} and
\eqref{m:eq:outward-mass} yields
\[
 K_{\mathrm{exc}}<C M r_*^{-3}+o(M)
                  =\frac{C}{B_0}K_{\mathrm{exc}}+o(M).
\]
Choose $B_0>2C$, increasing it also to ensure the initial scale choice in
Section~\ref{m:sec:conditioning}.  Since $K_{\mathrm{exc}}/M\to\infty$, the last
inequality is impossible.  This proves the existence of one universal finite
constant bounding the excess per nucleus in every strictly bound sector.

For an arbitrary strictly bound sector, the largest strictly bound sector of
the same system has at least as many electrons, so the same bound applies.
For every integer $N>Z+CM$, strict adjacent binding is therefore impossible.
Sector monotonicity gives $E_N\le E_{N-1}$, and exclusion of the strict
inequality gives $E_N=E_{N-1}$.  Taking $M=1$ proves both atomic assertions
for every real nuclear charge $Z\ge1$.
\end{proof}
